% !TEX root = main.tex
% Section 03: Architecture.
% Source of truth: docs/ARCHITECTURE.md.

\chapter{Architecture}

\section{Layered model}

\begin{verbatim}
+------------------------------------------------------------+
| 5. Experience                                              |
|    Browser extension . Desktop shell . CLI . future UIs    |
+------------------------------------------------------------+
| 4. Augmentation                                            |
|    Agents . Apps . Add-ons . Workflows                     |
+------------------------------------------------------------+
| 3. Intelligence                                            |
|    Cognitive Core . Local LLMs . Cloud LLMs                |
+------------------------------------------------------------+
| 2. Cognitive Infrastructure                                |
|    Living Archive . Context . Model Fabric                 |
+------------------------------------------------------------+
| 1. Resonant Core                                           |
|    Capability engine . Security . IPC . Host OS            |
+------------------------------------------------------------+
\end{verbatim}

Layer 1 is deterministic software. Layer 3 makes the system
intelligent. Layers 2 and 4 are policy and orchestration glue.

\section{Resonant Core (Layer 1)}

The Core is the single substrate that all clients (browser, desktop,
CLI, future UIs) speak to. It is the architectural center of the
system, not the browser.

Responsibilities:
\begin{itemize}
  \item Enforce every capability decision through deterministic code.
  \item Maintain the local user state root (\texttt{\textasciitilde/ResonantOS\_User} by default).
  \item Audit every privileged action.
  \item Hold provider credentials, secrets, and bridge tokens.
  \item Never depend on any AI component to make a security decision.
\end{itemize}

Initial implementation: Rust (per \adr{0001}).

\section{Cognitive Infrastructure (Layer 2)}

\begin{itemize}
  \item \textbf{Living Archive}: governed persistent memory.
  \item \textbf{Context}: working context assembly for the
        Cognitive Core.
  \item \textbf{Model Fabric}: unified router over local and remote
        intelligence.
\end{itemize}

The Model Fabric treats every model as a provider with measured
properties:
\begin{longtable}{p{2.4cm}p{10.0cm}}
\toprule
\textbf{Property} & \textbf{Meaning} \\
\midrule
latency & round-trip time for a representative request \\
cost & currency per million tokens (or per request for embedded models) \\
privacy & whether the request leaves the host \\
capability & measured quality on the lab's task suite \\
memory footprint & peak resident memory during inference \\
availability & probability the provider is reachable at request time \\
\bottomrule
\end{longtable}

\section{Intelligence (Layer 3)}

\begin{itemize}
  \item \textbf{Cognitive Core}: always-on small local model.
  \item \textbf{Local LLMs}: larger local specialists (optional).
  \item \textbf{Remote LLMs}: cloud / corporate / federated models
        (optional).
\end{itemize}

The Cognitive Core handles structured requests like:

\begin{lstlisting}[language=rccosjson,caption={Example request to the Cognitive Core.}]
{
  "task": "route",
  "input": "Open the budget spreadsheet and summarize it",
  "available_capabilities": ["filesystem.read", "spreadsheet.inspect"],
  "available_models": ["core-local", "local-7b", "cloud-gpt"],
  "context_summary": "..."
}
\end{lstlisting}

and returns:

\begin{lstlisting}[language=rccosjson,caption={Example response from the Cognitive Core.}]
{
  "intent": "summarize_document",
  "capabilities": ["filesystem.read", "spreadsheet.inspect"],
  "model_tier": "local",
  "memory_queries": ["budget-2026", "spreadsheet:budget"]
}
\end{lstlisting}

The deterministic Core then validates and executes.

\section{Augmentation (Layer 4)}

\begin{itemize}
  \item \textbf{Agents}: long-running autonomous workflows.
  \item \textbf{Apps}: scoped applications.
  \item \textbf{Add-ons}: third-party extensions.
  \item \textbf{Workflows}: user-defined multi-step flows.
\end{itemize}

\section{Experience (Layer 5)}

\begin{itemize}
  \item \textbf{Browser extension}: Chrome MV3 (carried over from baseline).
  \item \textbf{Desktop}: Tauri or native shell (new in stage 4).
  \item \textbf{CLI}: terminal-oriented client.
  \item \textbf{Mobile}: future.
\end{itemize}

\section{Open architectural questions}

\begin{enumerate}
  \item How should the Cognitive Core be trained? ResonantOS-specific
        fine-tuning of an existing ternary SLM vs.\ an off-the-shelf
        SLM?
  \item How is hardware discovery structured --- Rust trait, JSON
        contract, both?
  \item Where does IPC terminate? Loopback HTTP is the ResonantOS
        default; is that the right long-term transport?
  \item What is the upgrade story for the Cognitive Core model
        without breaking capability contracts?
\end{enumerate}
